\chapter{Анализ современных подходов к верификации подсистемы памяти с учётом моделей памяти}\label{ch:ch1}

В первой главе проанализированы архитектурные подходы к проектированию подсистемы памяти в многопроцессорных вычислительных системах.
Также приведен обзор существующих моделей памяти, рассмотрено их влияние на выполнение параллельных программ и отражение в практических реализациях вычислительных систем.
Представлены описания существующих моделей памяти и их способы задания.
Приведён обзор влияния модели памяти на операции с памятью на примере конкретной архитектуры.
Описано отражение моделей памяти на опыт использования различных языков программирования.
Проведён анализ существующих общепринятых решений верификации подсистемы памяти с учётом моделей памяти.

Процесорные технологии развиваются, предлагая всё более мощные и эффективные вычислительные системы.
По закону Мура~\cite{moore} рост числа транзисторов на кристалле удваивается каждые два года.
Но в последнее время рост числа транзисторов не коррелировал с ростом тактовой частоты, как было ранее~\cite{trend}.
На рисунке~\ref{fig:trend} показано изменение количества транзисторов, количества вычислительных ядер, тактовой частоты вычислительного ядра и энергопотребления.

\begin{figure}[ht]
    \centerfloat{
      \includegraphics[scale=2]{trend.jpg}
    }
    \caption{Изменение количества транзисторов, вычислительных ядер, тактовой частоты и энергопотребления.}
    \label{fig:trend}
\end{figure}

Рост тактовой частоты, который ранее являлся основным драйвером повышения производительности, замедлился.
По пути повышения быстродействия вычислительное ядро приблизилось к своим физическим и технологическим пределам~\cite{physLimits}.
Для дальнейшего роста производительности, инженеры обратились к другим стратегиям, и ключевым направлением развития стала многоядерность, что позволило выполнять несколько задач одновременно, повышая общую производительность.
Вследствие этого, сложность систем значительно увеличилась.
В них применяются многоуровневые иерархические архитектуры, параллелизм исполнения требует сложных механизмов синхронизаций, усложняя взаимодействие между различными компонентами~\cite{hieracy}.
Усложнение протоколов синхронизации подсистемы памяти сделало задачу верификации современных вычислительных систем актуальной.
Анализ современных подходов в решении проблемы верификации подсистемы памяти позволил выделить основные направления и методы решения задачи диссертационного исследования.

\section{Многопроцессорные системы}\label{sec:ch1/sec1}

Начнём рассмотрение многопроцессорных систем с изучения однопроцессорной.
Первым коммерчески доступным однокристальным микропроцессором был чип Intel 4004~\cite{intel4004}.
В чипе использовалось 16 регистров и прямое обращение к 4КБ ПЗУ и 640Б ОЗУ.
Схематично работа с памятью в Intel 4004 показана на рисунке~\ref{fig:intel4004}.

\begin{figure}[ht]
    \centerfloat{
      \includegraphics[scale=0.6]{intel4004-memory.png}
    }
    \caption{Схематичный вид работы с памятью в Intel 4004.}
    \label{fig:intel4004}
\end{figure}

Процессоры всегда работали быстрее, чем память.
Увеличение количества транзисторов в кристале позволяло реализовывать конвейерную и суперскалярную архитектуру~\cite{pipeline_superscalar}, что ещё больше повышало скорость работы процессора.
Разработчки памяти обычно использовали новые технологии для повышения ёмкости.
Это привело к простаиванию процессора в ожидании данных, после того как он сделал запрос в память.
Однако, медленная память -- это экономическая проблема.
Есть возможность построить память, которая будет работать так же быстро, как и процессор, но при этом её необходимо размещать прямо на микросхему процессора.
На самом деле, память, размещённая непосредственно на процессоре всегда существовала.
Это регистровая память, скорость обращение к которой такая же как и скорость самого процессора.
Но размещение большей памяти в процессор делает его больше и дороже, и существуют ограничения в размерах процессора, который можно сконструировать.
Таким образом, приходится выбирать между быстрой небольшой памятью и медленной памятью большего размера.

Существуют технологии, которые сочетают медленную большую память и быструю маленькую память.
Это позволяет получить и быструю память, и большую ёмкость.
Маленькая память с высокой скоростью работы называется кэш-память~\cite{cache}.
Основная идея кэш-памяти заключается в том, что в ней находятся данные, которые чаще всего используются.
Если процессору нажны данные, то вначале он их ищет в кэш-памяти.
И только если данных там нет, обращается в основную память.

Большинство программ содержат циклы, то есть блоки повторяющегося кода.
Благодаря тому, что выполняется ограниченный набор команд, соблюдаются принципы локальности~\cite{spatio-temp}.
Временной локальности -- если данные были использованны, то скорее всего они будут использованны ещё. И пространственной локальности -- если данные были использованны, то скорее всего будут использованы соседние данные.
Сама концепция кэш-памяти напрямую использует принцип временной локальности.
Сохраняя используемые данные для быстрого доступа, процессору нет необходимости обращаться в общую память для доступа к ним.
Принцип пространственной локальности уже используют алгоритмы вытеснения~\cite{cache-alg} и предзагрузки~\cite{prefetch}\fixme{find more apropreate cite} в кэш-память, исследование которых выходит за рамки данной работы.

Таком образом, схематичное взаимодействие процессора с памятью показано на рисунке~\ref{fig:memory_with_cache}.

\begin{figure}[ht]
    \centerfloat{
      \includegraphics[scale=0.5]{cpu-with-cache.png}
    }
    \caption{Схематичный вид взаимодействия процессора памятью.}
    \label{fig:memory_with_cache}
\end{figure}












Предметом рассмотрения многопроцессорных систем в работе является SMP(symmetric multiprocessing) архитектура~\cite{smp}, схематично изображённая на рисунке~\ref{fig:smp_arch}.

\begin{figure}[ht]
    \centerfloat{
      \includegraphics[scale=1]{smp-arch.png}
    }
    \caption{Схематичное вид SMP архитектуры.}
    \fixme{from smp}
    \label{fig:smp_arch}
\end{figure}

Главной особенностью данной архитектуры является наличие общей визической памяти, разделяемой всеми процессорами.
